\documentclass[11pt,a4paper]{article}

\usepackage[T1]{fontenc}
\usepackage[utf8]{inputenc}
\usepackage{lmodern}
\usepackage[margin=2.2cm]{geometry}
\usepackage{amssymb}
\usepackage{booktabs}
\usepackage{longtable}
\usepackage{array}
\usepackage[table]{xcolor}
\usepackage{enumitem}
\usepackage[hidelinks]{hyperref}

\definecolor{cut}{HTML}{EEF2F7}
\setlist{itemsep=2pt,topsep=4pt}
\setlength{\parindent}{0pt}
\setlength{\parskip}{6pt}
\renewcommand{\arraystretch}{1.15}
\newcolumntype{L}[1]{>{\raggedright\arraybackslash}p{#1}}

\title{Marathon Training Plan: Oct 2026 -- Aug 2027}
\date{Race: Saturday 28 August 2027}
\author{}

\begin{document}
\maketitle
\vspace{-1.5em}

% ------------------------------------------------------------------
\section*{Where you are now}

Your fitness is in decent shape; your consistency and load management are what need rebuilding. Recent short runs at 5:17--5:47\,/km are close to your pre-injury paces, but the last 9 weeks averaged under 10\,km per week with several 1--3 week gaps.

\begin{center}
\begin{tabular}{@{}L{3.1cm}L{2.9cm}L{1.9cm}L{1.6cm}L{4.4cm}@{}}
\toprule
\textbf{Period} & \textbf{Weekly km} & \textbf{Longest run} & \textbf{Runs/ week} & \textbf{Notes} \\
\midrule
28 Sep -- 4 Oct & 14.7 ($\approx$20 with Saturday's 5\,km) & 8.7 km & 2--3 & Two days in a row, no reported pain \\
27 Jul -- 27 Sep & 0--16 & 10.2 km & 0--3 & On-and-off comeback with gaps of 1--3 weeks \\
15 Jun -- 26 Jul & 0 & -- & 0 & Hip flexor break, about 6 weeks \\
18 May -- 14 Jun & 39--48 & 24 km & 3--4 & Marathon build, ended by the injury \\
\bottomrule
\end{tabular}
\end{center}

Three things in the pre-injury block likely set up the overuse injury:
\begin{itemize}
  \item \textbf{The long run grew too fast.} It went 18 $\rightarrow$ 20 $\rightarrow$ 22 $\rightarrow$ 24\,km in under three weeks, with only 5 days between the 22 and the 24.
  \item \textbf{The long run was too big a share of the week.} The 24\,km was 61\% of a 39\,km week; a safer ceiling is about 35--40\%.
  \item \textbf{Most runs were moderately hard.} Nearly every non-long run sat at 160--168\,bpm. Your truly easy runs (6:29\,/km at 147\,bpm, 5:58\,/km at 152\,bpm) were rare.
\end{itemize}

You also ran 5\,km and 10\,km in the four days after the 24\,km, which may have prolonged the problem. None of this is unusual; it is the most common pattern behind first-marathon injuries.

% ------------------------------------------------------------------
\section*{Is the timeline realistic?}

Yes. Your race on \textbf{Saturday 28 August 2027} is 47 weeks away, which is still plenty. A runner with two half marathons behind them typically needs 16--20 weeks of marathon-specific training; this plan keeps a full 18-week marathon block (weeks 30--47) and shortens only the Base and Build phases. The 8-week Rebuild for your hip is unchanged. The real risk is not time but repeating the overload that hurt your hip.

The plan peaks at about 66\,km per week with three long runs of 30--32\,km. Most of the marathon block falls in summer, so see the note on hot days in the sessions section.

\subsection*{The rules that make the plan work}
\begin{enumerate}
  \item \textbf{Easy means easy.} Run easy days at 6:00--6:40\,/km, or by heart rate at roughly 150\,bpm or below (your data shows 147--155 at those paces). If heart rate climbs, slow down. About 80\% of weekly running should be at this effort.
  \item \textbf{One variable at a time.} Weekly km rises by no more than about 10\%, and every fourth week is a cutback of 20--25\%.
  \item \textbf{Long run caps.} It grows by at most 1--2\,km per week and stays at or under 40\% of weekly km until the marathon block. There the 28--32\,km runs reach 45--50\% of the week, so they come only every other week with lighter long runs between.
  \item \textbf{Hip pain rule.} Mild discomfort up to 2/10 that does not worsen during the run and is gone by the next morning is fine. Anything above 3/10, pain that changes your stride, or pain still there the next morning means stop, rest 2--3 days, and restart one week lower. If it recurs twice, see a physiotherapist.
  \item \textbf{Strength twice a week.} Twenty minutes of hip and core work (section below) is part of the plan, not an extra.
  \item \textbf{Rest after long runs.} The day after a long run is rest or very easy; no hard session within 48 hours.
\end{enumerate}

I am not a medical professional; if the hip keeps nagging, a physio assessment early on is worth more than any plan.

% ------------------------------------------------------------------
\section*{Phases at a glance}

Each phase ends with a simple check. Don't pass it on the calendar alone: if the check isn't met, repeat the previous week until it is.

{\small
\setlength{\tabcolsep}{4pt}
\begin{center}
\begin{tabular}{@{}L{1.7cm}L{2.5cm}L{1.6cm}L{1.6cm}L{4.0cm}L{3.7cm}@{}}
\toprule
\textbf{Phase} & \textbf{Weeks / dates} & \textbf{Km per week} & \textbf{Long run} & \textbf{Focus} & \textbf{Check before moving on} \\
\midrule
Rebuild & Wk 1--8 \newline 5 Oct -- 29 Nov & 20 $\rightarrow$ 31 & 8 $\rightarrow$ 12 km & Easy running only, 3 runs a week; start strength work twice a week & 3 pain-free weeks in a row \\
\addlinespace
Base & Wk 9--20 \newline 30 Nov -- 21 Feb & 32 $\rightarrow$ 46 & 13 $\rightarrow$ 18 km & 4 runs a week, strides, hill sprints; fartlek and steady runs from late Jan & 45\,km weeks feel routine and the hip stays quiet \\
\addlinespace
Build & Wk 21--29 \newline 22 Feb -- 25 Apr & 47 $\rightarrow$ 55 & 18 $\rightarrow$ 21 km & One threshold or tempo session a week; half marathon race in week 28 & Half marathon result sets your marathon pace \\
\addlinespace
Marathon block & Wk 30--44 \newline 26 Apr -- 8 Aug & 50 $\rightarrow$ 66 & 20 $\rightarrow$ 32 km & 5 runs a week, MP inside long runs; three long runs of 30--32 km & Last 32\,km long run done \\
\addlinespace
Taper & Wk 45--47 \newline 9 -- 28 Aug & 50 $\rightarrow$ 15 & 22 $\rightarrow$ 16 km & Volume drops, short MP pieces stay; rest, sleep, rehearse race-day kit & \textbf{Race day: Saturday 28 August 2027, 42.2\,km} \\
\bottomrule
\end{tabular}
\end{center}
}

% ------------------------------------------------------------------
\section*{Week by week}

Tick the week off when done. Km and long-run figures are targets, not minimums; landing 10\% under is fine. MP = goal marathon pace, set from your half marathon in week 28. Cutback weeks are shaded. Session details are in the next section.

{\small
\setlength{\tabcolsep}{4pt}
\begin{longtable}{@{}c r L{2.3cm} L{1.6cm} L{1.2cm} L{1.1cm} L{6.3cm}@{}}
\toprule
\textbf{\checkmark} & \textbf{Wk} & \textbf{Dates} & \textbf{Phase} & \textbf{Km} & \textbf{Long run} & \textbf{Key session / notes} \\
\midrule
\endfirsthead
\toprule
\textbf{\checkmark} & \textbf{Wk} & \textbf{Dates} & \textbf{Phase} & \textbf{Km} & \textbf{Long run} & \textbf{Key session / notes} \\
\midrule
\endhead
\bottomrule
\endfoot
$\square$ & 1 & 5--11 Oct & Rebuild & 20 & 8 & 3 easy runs (5, 7, 8 km) \\
$\square$ & 2 & 12--18 Oct & Rebuild & 22 & 9 & 3 easy runs \\
$\square$ & 3 & 19--25 Oct & Rebuild & 24 & 10 & 3 easy runs \\
\rowcolor{cut}$\square$ & 4 & 26 Oct--1 Nov & Rebuild & 19 & 8 & Cutback \\
$\square$ & 5 & 2--8 Nov & Rebuild & 26 & 11 & Add 4 $\times$ 20\,s strides after one easy run \\
$\square$ & 6 & 9--15 Nov & Rebuild & 28 & 12 & Optional 4th run of 4--5 km \\
$\square$ & 7 & 16--22 Nov & Rebuild & 31 & 12 & 4 runs \\
\rowcolor{cut}$\square$ & 8 & 23--29 Nov & Rebuild & 25 & 10 & Cutback. Gate: 3 pain-free weeks before moving on \\
\midrule
$\square$ & 9 & 30 Nov--6 Dec & Base & 32 & 13 & Strides twice a week from now on \\
$\square$ & 10 & 7--13 Dec & Base & 34 & 13 &  \\
$\square$ & 11 & 14--20 Dec & Base & 36 & 14 &  \\
\rowcolor{cut}$\square$ & 12 & 21--27 Dec & Base & 28 & 11 & Cutback \\
$\square$ & 13 & 28 Dec--3 Jan & Base & 37 & 15 & Add 6 $\times$ 10\,s hill sprints to one easy run \\
$\square$ & 14 & 4--10 Jan & Base & 39 & 15 &  \\
$\square$ & 15 & 11--17 Jan & Base & 41 & 16 &  \\
\rowcolor{cut}$\square$ & 16 & 18--24 Jan & Base & 32 & 12 & Cutback \\
$\square$ & 17 & 25--31 Jan & Base & 43 & 16 & Fartlek: 8 $\times$ 1 min steady / 1 min easy \\
$\square$ & 18 & 1--7 Feb & Base & 45 & 17 & Fartlek: 6 $\times$ 2 min / 1 min \\
$\square$ & 19 & 8--14 Feb & Base & 46 & 18 & Steady 20 min \\
\rowcolor{cut}$\square$ & 20 & 15--21 Feb & Base & 37 & 14 & Cutback \\
\midrule
$\square$ & 21 & 22--28 Feb & Build & 47 & 18 & Threshold 3 $\times$ 8 min \\
$\square$ & 22 & 1--7 Mar & Build & 49 & 19 & Threshold 3 $\times$ 10 min \\
$\square$ & 23 & 8--14 Mar & Build & 51 & 20 & Tempo 20 min \\
\rowcolor{cut}$\square$ & 24 & 15--21 Mar & Build & 41 & 15 & Cutback; 5 km effort test to update paces \\
$\square$ & 25 & 22--28 Mar & Build & 53 & 20 & Threshold 4 $\times$ 8 min \\
$\square$ & 26 & 29 Mar--4 Apr & Build & 55 & 21 & Tempo 25 min; last 5 km of long run at half-marathon effort \\
$\square$ & 27 & 5--11 Apr & Build & 47 & 15 & Mini-taper: 2 $\times$ 10 min threshold \\
$\square$ & 28 & 12--18 Apr & Build & 42 & 21.1 & \textbf{Half marathon race} (checkpoint, sets MP) \\
$\square$ & 29 & 19--25 Apr & Build & 36 & 14 & Recovery, easy only \\
\midrule
$\square$ & 30 & 26 Apr--2 May & Marathon & 50 & 20 & Threshold 3 $\times$ 10 min \\
$\square$ & 31 & 3--9 May & Marathon & 54 & 22 & Tempo 25 min; medium-long run 14 km \\
$\square$ & 32 & 10--16 May & Marathon & 57 & 24 & Long run with last 6 km at MP \\
\rowcolor{cut}$\square$ & 33 & 17--23 May & Marathon & 46 & 18 & Cutback \\
$\square$ & 34 & 24--30 May & Marathon & 58 & 25 & Threshold 4 $\times$ 10 min; medium-long 15 km \\
$\square$ & 35 & 31 May--6 Jun & Marathon & 60 & 26 & Long run with 2 $\times$ 5 km at MP \\
$\square$ & 36 & 7--13 Jun & Marathon & 62 & 28 & Midweek 10 km at MP (inside a 14 km run) \\
\rowcolor{cut}$\square$ & 37 & 14--20 Jun & Marathon & 50 & 20 & Cutback \\
$\square$ & 38 & 21--27 Jun & Marathon & 62 & 28 & Long run with 10 km at MP in the middle \\
$\square$ & 39 & 28 Jun--4 Jul & Marathon & 64 & 30 & Easy long run; threshold 3 $\times$ 12 min midweek \\
$\square$ & 40 & 5--11 Jul & Marathon & 64 & 26 & Long run with 14 km at MP \\
\rowcolor{cut}$\square$ & 41 & 12--18 Jul & Marathon & 52 & 20 & Cutback \\
$\square$ & 42 & 19--25 Jul & Marathon & 66 & 32 & Easy long run; practise race fuelling \\
$\square$ & 43 & 26 Jul--1 Aug & Marathon & 60 & 24 & Long run with 16 km at MP (dress rehearsal) \\
$\square$ & 44 & 2--8 Aug & Marathon & 66 & 32 & Last long run, easy \\
\midrule
$\square$ & 45 & 9--15 Aug & Taper & 50 & 22 & Long run with 8 km at MP \\
$\square$ & 46 & 16--22 Aug & Taper & 38 & 16 & 2 $\times$ 3 km at MP \\
$\square$ & 47 & 23--29 Aug & Taper & 15 + race & 42.2 & 3 short easy runs with strides Mon--Thu, rest Fri; \textbf{race Saturday 28 Aug} \\
\end{longtable}
}

The week of 28 Sep -- 4 Oct (now) counts as week 0: your Wednesday, Thursday and Saturday runs make about 20\,km, which is exactly where week 1 starts. If any of the gaps in August and September were the hip flaring up, repeat weeks 1--2 at 3 $\times$ 5\,km before progressing.

% ------------------------------------------------------------------
\section*{Weekly templates and sessions}

Days are movable; you have done long runs on Mondays and Wednesdays before, and that is fine. Keep the order: no two hard days (long run or quality session) back to back. Since the race is on a Saturday, try to do at least the last few long runs on a Saturday morning to rehearse the routine.

\begin{center}
\begin{tabular}{@{}L{2.2cm}L{1.3cm}L{1.9cm}L{9.6cm}@{}}
\toprule
\textbf{Phase} & \textbf{Weeks} & \textbf{Runs} & \textbf{Typical week} \\
\midrule
Rebuild & 1--8 & 3 (4 from wk 6) & Easy $\cdot$ rest $\cdot$ easy $\cdot$ rest $\cdot$ long $\cdot$ rest $\cdot$ rest \\
Base & 9--20 & 4 & Easy + strides $\cdot$ rest $\cdot$ easy $\cdot$ rest $\cdot$ easy + strides/fartlek $\cdot$ long $\cdot$ rest \\
Build & 21--29 & 4--5 & Easy $\cdot$ quality $\cdot$ rest $\cdot$ easy $\cdot$ rest $\cdot$ long $\cdot$ optional easy 5\,km \\
Marathon & 30--44 & 5 & Easy $\cdot$ quality $\cdot$ easy $\cdot$ medium-long $\cdot$ rest $\cdot$ long $\cdot$ rest \\
Taper & 45--47 & 4--5 & Same shape as Marathon, shorter, with short MP pieces \\
\bottomrule
\end{tabular}
\end{center}

Strength (next section) goes on two easy or rest days, not the day before the long run.

\begin{center}
\begin{tabular}{@{}L{3.2cm}L{12.0cm}@{}}
\toprule
\textbf{Session} & \textbf{How to run it} \\
\midrule
Easy run & 6:00--6:40\,/km, about 150\,bpm or below, fully conversational. \\
Long run & Same effort as easy. Start slow; the last third may drift slightly faster. From 20\,km, carry water and practise gels or sports drink. \\
Strides & 4--6 $\times$ 20\,s smooth and fast (not sprinting) on flat ground, walk back to recover. After an easy run. \\
Hill sprints & 6--8 $\times$ 10\,s hard uphill, walk down fully recovered. Builds strength with little impact. \\
Fartlek / steady & Comfortably hard but in control, around 5:30--5:45\,/km for you today. \\
Threshold / tempo & ``Comfortably hard'': you could say a few words but not chat, roughly 160--168\,bpm for you. Intervals take 2 min easy jog between. 2\,km easy warm-up and cool-down. \\
Medium-long run & 14--18\,km easy, midweek, in the marathon block. \\
MP segments & At goal marathon pace inside a longer run. Set MP from your week-28 half: half-marathon pace + about 4--6\% (for example, 5:20\,/km half pace $\rightarrow$ 5:33--5:39\,/km MP). \\
5 km effort test & Warm up 2\,km, then a hard but even 5\,km. Use it to sanity-check threshold pace. \\
Hot days & Common in the June--August marathon block. Run easy days by heart rate rather than pace, do long runs early in the morning, drink to thirst plus electrolytes on runs over 90 minutes, and let MP segments run 5--10\,s/km slower when it is above about 20\,$^{\circ}$C. \\
\bottomrule
\end{tabular}
\end{center}

% ------------------------------------------------------------------
\section*{Strength and hip routine}

Twice a week, about 20 minutes, from week 1. Strong glutes and hip flexors share the load your hip flexor took alone during the long runs. Start with bodyweight; add a band or weight once 3 sets feel easy.

\begin{center}
\begin{tabular}{@{}L{7.4cm}L{3.9cm}L{3.9cm}@{}}
\toprule
\textbf{Exercise} & \textbf{Sets $\times$ reps} & \textbf{Targets} \\
\midrule
Glute bridge $\rightarrow$ single-leg bridge & 3 $\times$ 12 (each leg) & Glutes, hamstrings \\
Standing banded knee drive (or lying straight-leg raise) & 3 $\times$ 12 each side, slow lowering & Hip flexors \\
Side plank with top-leg lift & 3 $\times$ 20--30\,s each side & Hip abductors, core \\
Reverse lunge or split squat & 3 $\times$ 10 each leg & Quads, glutes, balance \\
Single-leg calf raise & 3 $\times$ 15 each leg & Calves, Achilles \\
Dead bug & 3 $\times$ 10 each side & Deep core, pelvic control \\
\bottomrule
\end{tabular}
\end{center}

After runs, a gentle kneeling hip flexor stretch (2 $\times$ 30\,s per side) is fine if it feels good; skip it if it provokes the old pain.

% ------------------------------------------------------------------
\section*{Adjusting the plan}

\begin{itemize}
  \item \textbf{No half marathon in mid-April.} Race one in early May instead and swap it for the week-31 long run, with an easy recovery week after. If there is none, run a solo 18\,km with the last 10\,km at half-marathon effort and set MP from that.
  \item \textbf{Missed up to a week.} Carry on with the plan as written.
  \item \textbf{Missed 1--2 weeks.} Go back one week per missed week, then continue. Before week 30, the time comes out of Base or Build; never shorten the marathon block or the taper.
  \item \textbf{Missed more than 2 weeks, or illness.} Restart at about 60\% of your last full week and rebuild over 2--3 weeks.
  \item \textbf{Hip pain returns.} Follow rule 4. If you need more than 3 days off, also drop quality sessions for the following week.
  \item \textbf{Feeling great.} Stay with the plan. Spare energy is better spent on strength work and sleep than extra kilometres.
\end{itemize}

\end{document}
